% !TeX root=../main.tex

\chapter{مقدمه و بیان مسئله}

\section{مقدمه}
\label{sec:1-1}

در سلف‌های دانشگاه تهران غذا بر مبنای تعداد رزروهای ثبت‌شده پخته می‌شود، ولی بخشی از این
رزروها هرگز تحویل گرفته نمی‌شوند و غذای متناظرشان، که هزینه‌اش پرداخت شده، مصرف‌کننده‌ای
ندارد. کاهش این هدررفت با پخت کمتر، بدون تخمین اینکه چقدر کمتر، ریسک کمبود غذا را می‌سازد
که هزینه‌اش برای دانشگاه سنگین‌تر است. مسئله بنابراین یک مسئله‌ی پیش‌بینی است و نه صرفاً یک
مسئله‌ی عملیاتی: پیش از هر وعده باید تخمین زد چه نسبتی از رزروها دریافت نخواهد شد. این
پژوهش همان نسبت را با رگرسیون کوانتایل روی سری‌های زمانی چندگانه برآورد می‌کند و آن را به
عددی برای پخت تبدیل می‌کند که هزینه‌ی نابرابر کمبود و مازاد در آن لحاظ شده است.

\section{تاریخچه‌ای از موضوع تحقیق}
\label{sec:1-2}

تخصیص غذا در دانشگاه تهران از طریق یک سامانه‌ی رزرو اینترنتی انجام می‌شود. دانشجو غذای هر
وعده را دست‌کم ۷۲ ساعت پیش از سرو رزرو می‌کند و پس از بسته‌شدن مهلت امکان لغو ندارد؛ پس در
لحظه‌ای که آشپزخانه باید درباره‌ی مقدار پخت تصمیم بگیرد، تعداد رزرو عددی قطعی و کاملاً معلوم
است. هر رزرو به ترکیبی از روز، وعده، سلف و نوع غذا نسبت داده می‌شود و هنگام سرو، تحویل با
کارت دانشجویی یا با کد ثبت می‌شود. رویه‌ی کنونی مدیریت پخت همین عدد رزرو را مبنا می‌گیرد؛
مصاحبه‌های آغاز پروژه نشان داد پیمانکاران در عمل گاهی بر پایه‌ی تجربه‌ی شخصی کمتر از آن
می‌پزند، بی‌آنکه این کاهش با کارفرما هماهنگ شود یا در صورتحساب بازتاب یابد. تصمیم واقعی
پخت بنابراین بیرون از سازوکار رسمی و بدون مبنای مستند گرفته می‌شود.

\section{شرح مسئله تحقیق}
\label{sec:1-3}

داده‌ی سامانه‌ی تغذیه در بازه‌ی آذر ۱۴۰۲ تا خرداد ۱۴۰۳، شامل ۱۴۲ روز سرو در ۳۰ سلف و شش
شهر، نشان می‌دهد $8.05\%$ کل رزروهای ثبت‌شده تحویل گرفته نشده‌اند. این عدد نسبت مجموع
رزروهای دریافت‌نشده به مجموع کل رزروهاست؛ میانگین ساده‌ی نرخ سلول‌ها $9.59\%$ است، یعنی
حدود ۱۹ درصد بالاتر، چون به سلف‌های کوچک و پرنوسان وزنی بیش از سهمشان در حجم غذا می‌دهد.
حجم مطلق هدررفت بیش از ۱۶۲ هزار پرس در پنج ماه است که با قیمت تقریبی ۱۲۰ هزار تومان برای
هر پرس، نزدیک به ۱۹ میلیارد تومان در همین بازه و سالانه‌شده حدود ۵۰ میلیارد تومان می‌شود.

هزینه‌ی خطا میان سه گروه درگیر یکسان توزیع نمی‌شود: مدیریت سلف مقدار پخت را تعیین می‌کند،
پیمانکار خرید و پخت را بر همان عدد می‌بندد، و دانشجو که در تصمیم نقشی ندارد پیامد خطای رو
به پایین را به شکل کمبود غذا تجربه می‌کند. پخت بیش از نیاز هزینه‌ی مالی دارد و قابل جبران
است، ولی پخت کمتر از نیاز به قطع خدمت می‌انجامد. همین عدم‌تقارن، نه دقت متقارن پیش‌بینی،
شکل تابع هدف این پژوهش را تعیین می‌کند.

\section{تعریف موضوع تحقیق}
\label{sec:1-4}

واحد تحلیل چهارتایی $(d,m,r,f)$ است: روز، وعده، سلف و نوع غذا، یعنی همان سطحی که آشپزخانه
برای آن یک عدد پخت می‌گیرد. کمیت پیش‌بینی‌شونده نرخ عدم‌دریافت $\rho$ است، نسبت رزروهای
دریافت‌نشده به کل رزروهای همان سلول، و آنچه برآورد می‌شود میانگین شرطی این نرخ نیست بلکه
کوانتایل $\tau$اُم توزیع شرطی آن است؛ دلیل هر دو انتخاب در بندهای \ref{sec:2-2-3} و
\ref{sec:3-2-1} آمده است. همین کمیت را می‌توان در سطوح مختلفی از تجمیع مدل کرد، از سطح
سلول تا سطح تک‌رزرو هر دانشجو، و این پژوهش سطح داده را خودش یک محور آزمایش قرار داد تا اثر
الگوریتم و اثر سطح تجمیع در مقایسه‌ها آمیخته نشوند. قید حاکم بر کل کار، قاعده‌ی برش
اطلاعاتی است: هر پیش‌بینی تنها با اطلاعاتی ساخته می‌شود که در لحظه‌ی واقعی تصمیم در دسترس
بوده‌اند. این لحظه برای ناهار و شام یکی نیست و تفاوتشان پیامدی دارد که در نگاه اول دیده
نمی‌شود، چنان‌که بند \ref{sec:3-2-2} نشان می‌دهد.

\section{اهداف و آرمان‌های کلی تحقیق}
\label{sec:1-5}

هدف پژوهش برآورد مقدار قابل‌پخت هر وعده در سطح $(d,m,r,f)$ است، به‌گونه‌ای که هدررفت مالی
نسبت به رویه‌ی کنونی کاهش یابد بدون آنکه ریسک کمبود غذا بالا برود، و سنجش آن با زیان
\lr{Pinball} در برابر خط پایه‌ای انجام می‌شود که خودِ پروژه می‌سازد، چون هیچ خط پایه‌ی
معتبری از پیش موجود نبود. انتخاب این معیار به‌جای \lr{RMSE} از همان عدم‌تقارن هزینه می‌آید:
\lr{RMSE} با میانگین شرطی کمینه می‌شود و جهت خطا برایش تفاوتی ندارد، در حالی که اینجا
کم‌برآوردکردن نرخ عدم‌دریافت به هدررفت و بیش‌برآوردکردنش به کمبود غذا می‌انجامد. \lr{RMSE}،
\lr{MAE} و ضریب تعیین نگه داشته شده‌اند، اما تنها برای مقایسه با ادبیات موجود و نه به‌عنوان
مبنای انتخاب مدل.

محدوده‌ی کار با چند فرض صریح بسته شده است: منوی هر وعده در لحظه‌ی رزرو معلوم است، تعداد رزرو
پس از بسته‌شدن مهلت تغییر نمی‌کند، پخت تک‌مرحله‌ای است و امکان پخت تکمیلی پس از دیدن تقاضا
وجود ندارد، ظرفیت آشپزخانه قید فعالی نیست، و سرنوشت غذای مازاد و سیاست جریمه‌ی عدم‌دریافت
بیرون از محدوده‌ی پژوهش‌اند. اثرگذارترین محدودیت اما پوشش زمانی داده است: پنج ماه آذر تا
خرداد برای الگوی هفتگی کافی است ولی هیچ اطلاعاتی درباره‌ی تابستان، شروع ترم و دوره‌ی
امتحانات ندارد، پس در این گزارش هیچ ادعای تعمیمی به آن بازه‌ها مطرح نمی‌شود.

\section{روش انجام تحقیق}
\label{sec:1-6}

کار در یازده فاز اجرا شد که شکل \ref{fig:phase-flow} ترتیب و وابستگی آن‌ها را نشان می‌دهد.
مسیر اصلی از پاکسازی و یکپارچه‌سازی داده‌ی خام آغاز می‌شود، از کاوش داده و مهندسی ویژگی
می‌گذرد، پس از تثبیت پروتکل اعتبارسنجی زمانی وارد مدل‌سازی می‌شود، و با ارزیابی،
تفسیرپذیری و ساخت لایه‌ی تصمیم پایان می‌یابد. این توالی خطی نیست: یافته‌های ارزیابی
می‌توانستند فهرست ویژگی‌ها را تغییر دهند و بازگشت به فازهای پیشین بخشی از طراحی بود، به
شرطی که مجموعه‌ی آزمون قفل‌شده در این رفت‌وبرگشت‌ها دست‌نخورده بماند.

میان فازها پنج دروازه‌ی کیفیت قرار گرفت که جدول \ref{tab:gates} اثر عملی هر کدام را
می‌آورد. هر دروازه یک ریسک روش‌شناختی مشخص را پوشش می‌داد: پیش‌رفتن با داده‌ای که هنوز
آزمون‌های سازگاری‌اش را پاس نکرده، ساختن ویژگی پیش از شناختن ساختار داده، آموزش مدل پیچیده
بدون مرجعی که بتوان برتری را نسبت به آن سنجید، و باز کردن زودهنگام مجموعه‌ی آزمون.

\begin{figure}[!ht]
\centering
\resizebox{\textwidth}{!}{%
\begin{tikzpicture}[
  font=\small,
  ph/.style={draw, rounded corners=2pt, minimum width=2.6cm, minimum height=1.05cm,
             align=center, fill=black!4},
  crit/.style={ph, fill=black!14, thick},
  gate/.style={draw, rounded corners=1pt, fill=black!70, text=white,
               font=\footnotesize\bfseries, inner sep=2pt},
  ar/.style={->, thick},
  fb/.style={->, dashed, gray}
]
% ردیف اول: راست به چپ
\node[ph]   (p0) at (11.4,5) {\rl{فاز ۰}\\\rl{زیرساخت}};
\node[ph]   (p1) at (7.6,5)  {\rl{فاز ۱}\\\rl{تعریف مسئله}};
\node[ph]   (p2) at (3.8,5)  {\rl{فاز ۲}\\\rl{جمع‌آوری داده}};
\node[crit] (p3) at (0,5)    {\rl{فاز ۳}\\\rl{پاکسازی}};
% ردیف دوم: چپ به راست
\node[crit] (p4) at (0,3)    {\rl{فاز ۴}\\\rl{کاوش داده}};
\node[crit] (p5) at (3.8,3)  {\rl{فاز ۵}\\\rl{مهندسی ویژگی}};
\node[ph]   (p6) at (7.6,3)  {\rl{فاز ۶}\\\rl{پروتکل ارزیابی}};
% ردیف سوم: راست به چپ
\node[crit] (p7) at (11.4,1) {\rl{فاز ۷}\\\rl{مدل‌سازی}};
\node[crit] (p8) at (7.6,1)  {\rl{فاز ۸}\\\rl{ارزیابی}};
\node[ph]   (p9) at (3.8,1)  {\rl{فاز ۹}\\\rl{تفسیرپذیری}};
\node[ph]   (p10) at (0,1)   {\rl{فاز ۱۰}\\\rl{لایه‌ی تصمیم}};
% ردیف چهارم
\node[ph, minimum width=8cm] (p11) at (5.7,-1.4) {\rl{فاز ۱۱ — مستندسازی و ارائه (موازی با همه‌ی فازها)}};

\draw[ar] (p0) -- (p1);
\draw[ar] (p1) -- (p2);
\draw[ar] (p2) -- (p3);
\draw[ar] (p3) -- (p4) node[gate, midway] {\lr{M1}};
\draw[ar] (p4) -- (p5) node[gate, midway] {\lr{M2}};
\draw[ar] (p5) -- (p6);
\draw[ar] (p6) -| (p7) node[gate, pos=0.42] {\lr{M3}};
\draw[ar] (p7) -- (p8) node[gate, midway] {\lr{M4}};
\draw[ar] (p8) -- (p9);
\draw[ar] (p9) -- (p10);
\draw[ar] (p10) |- (p11) node[gate, pos=0.30] {\lr{M5}};
\draw[fb] (p8) -- (p5);
\end{tikzpicture}}
\caption{جریان یازده فاز پروژه و پنج دروازه‌ی کیفیت. کادرهای پررنگ‌تر مسیر بحرانی‌اند؛
پیکان خط‌چین نمونه‌ای از مسیرهای بازگشت مجاز (ارزیابی $\rightarrow$ مهندسی ویژگی) است
که ماهیت حلقوی فرآیند را نشان می‌دهد.}
\label{fig:phase-flow}
\end{figure}

\begin{table}[!ht]
\centering
\singlespacing
\caption{دروازه‌های کیفیت و اثر عملی هر کدام}
\label{tab:gates}
\vspace{2mm}
\begin{tabular}{|p{1.3cm}|p{2.4cm}|p{9cm}|}
\hline
\textbf{دروازه} & \textbf{پایان فاز} & \textbf{چه چیزی را متوقف کرد} \\
\hline
\lr{M1} & ۳ (پاکسازی) & تا وقتی یک اسکریپت واحد داده‌ی خام را به مجموعه‌ی پردازش‌شده نرساند و آزمون‌های تقاطعی و هش \lr{snapshot} ثبت نشوند، هیچ ویژگی‌ای ساخته نشد \\
\hline
\lr{M2} & ۴ (کاوش داده) & پیش از پرشدن دفتر حقایق داده و تصمیم درباره‌ی سراسری یا خوشه‌ای بودن مدل، فاز ویژگی شروع نشد \\
\hline
\lr{M3} & ۶ (پروتکل ارزیابی) & مؤثرترین دروازه: تا پیش از اجرای همه‌ی خط پایه‌ها و تثبیت پروتکل اعتبارسنجی زمانی، هیچ مدل پیچیده‌ای آموزش ندید؛ بدون مرجع مقایسه، برتری هر مدل تنها با معیاری سرخود قابل‌ادعا بود \\
\hline
\lr{M4} & ۷ (مدل‌سازی) & هر پیکربندی وضعیت مستندی گرفت («تکمیل‌شده» یا «ناسازگار-مستند») و هر اجرا با هش \lr{fold} یا \lr{snapshot} نامنطبق از مقایسه کنار رفت؛ پنجره‌ی انتهایی در کل فاز دست‌نخورده ماند \\
\hline
\lr{M5} & ۸ تا ۱۰ و گزارش & مجموعه‌ی آزمون فقط یک بار باز شد؛ پس از آن هیچ تغییری در مدل انجام نشد \\
\hline
\end{tabular}
\end{table}

\section{ساختار پایان‌نامه}
\label{sec:1-7}

فصل دوم مفاهیم لازم برای خواندن بقیه‌ی گزارش را می‌آورد: مسئله‌ی روزنامه‌فروش و زیان
نامتقارن، رگرسیون کوانتایل و پیش‌بینی احتمالاتی، اعتبارسنجی و نشت در داده‌ی زمانی،
خانواده‌های مدل، ابزارهای آماری به‌کاررفته، و مروری بر پیشینه. فصل سوم روش پیشنهادی را
تشریح می‌کند، از صورت‌بندی مسئله و قاعده‌ی برش تا مهندسی ویژگی، غربالگری مدل‌ها،
کالیبراسیون و لایه‌ی تصمیم، و در کنار آن منابع داده، پاکسازی، یافته‌های کاوش و ممیزی نشت
را گزارش می‌کند. فصل چهارم همان روش را می‌سنجد: محیط اجرا، معیارها و پروتکل اعتبارسنجی،
خط پایه‌ها، نتایج روی مجموعه‌ی قفل‌شده، و تحلیل خطا، تفسیرپذیری و اثر کسب‌وکاری. فصل پنجم
جمع‌بندی، محدودیت‌ها و پیشنهادهای کار آینده را می‌آورد و فصل ششم فهرست مراجع است.
